
\documentclass{article}
\usepackage{colortbl}
\usepackage{makecell}
\usepackage{multirow}
\usepackage{supertabular}

\begin{document}

\newcounter{utterance}

\centering \large Interaction Transcript for game `hot\_air\_balloon', experiment `air\_balloon\_survival\_en\_complexity\_easy', episode 2 with gemini{-}3{-}flash{-}t1.0.
\vspace{24pt}

{ \footnotesize  \setcounter{utterance}{1}
\setlength{\tabcolsep}{0pt}
\begin{supertabular}{c@{$\;$}|p{.15\linewidth}@{}p{.15\linewidth}p{.15\linewidth}p{.15\linewidth}p{.15\linewidth}p{.15\linewidth}}
    \# & $\;$A & \multicolumn{4}{c}{Game Master} & $\;\:$B\\
    \hline

    \theutterance \stepcounter{utterance}  
    & & \multicolumn{4}{p{0.6\linewidth}}{
        \cellcolor[rgb]{0.9,0.9,0.9}{
            \makecell[{{p{\linewidth}}}]{
                \texttt{\tiny{[P1$\langle$GM]}}
                \texttt{You are participating in a collaborative negotiation game.} \\
\\ 
\texttt{Together with another participant, you must agree on a single set of items that will be kept. Each of you has your own view of how much each item matters to you (importance). You do not know how the other participant values the items. Additionally, you are given the effort each item demands.} \\
\texttt{You may only agree on a set if the total effort of the selected items does not exceed a shared limit:} \\
\\ 
\texttt{LIMIT = 7868} \\
\\ 
\texttt{Here are the individual item effort values:} \\
\\ 
\texttt{Item effort = \{"B21": 552, "B00": 560, "A08": 76, "A27": 950, "B94": 248, "B53": 871, "A23": 391, "C65": 143, "A89": 293, "B67": 208, "B45": 678, "C34": 737, "C12": 407, "A05": 366, "B60": 768, "C92": 859, "B69": 183, "A84": 231, "C06": 305, "B12": 727, "A16": 148, "B85": 357, "C35": 504, "B10": 549, "C15": 299, "B09": 91, "B16": 527, "C81": 848, "A58": 307, "C24": 214, "A30": 723, "A66": 475, "B75": 23, "B46": 298, "C70": 820\}} \\
\\ 
\texttt{Here is your personal view on the importance of each item:} \\
\\ 
\texttt{Item importance values = \{"B21": 138, "B00": 583, "A08": 868, "A27": 822, "B94": 783, "B53": 65, "A23": 262, "C65": 121, "A89": 508, "B67": 780, "B45": 461, "C34": 484, "C12": 668, "A05": 389, "B60": 808, "C92": 215, "B69": 97, "A84": 500, "C06": 30, "B12": 915, "A16": 856, "B85": 400, "C35": 444, "B10": 623, "C15": 781, "B09": 786, "B16": 3, "C81": 713, "A58": 457, "C24": 273, "A30": 739, "A66": 822, "B75": 235, "B46": 606, "C70": 105\}} \\
\\ 
\texttt{Goal:} \\
\\ 
\texttt{Your goal is to negotiate a shared set of items that benefits you as much as possible (i.e., maximizes total importance to YOU), while staying within the LIMIT. You are not required to make a PROPOSAL in every message {-} you can simply negotiate as well. All tactics are allowed!} \\
\\ 
\texttt{Interaction Protocol:} \\
\\ 
\texttt{You may only use the following structured formats in a message:} \\
\\ 
\texttt{PROPOSAL: \{'A', 'B', 'C', …\}} \\
\texttt{Propose keeping exactly those items.} \\
\texttt{REFUSE: \{'A', 'B', 'C', …\}} \\
\texttt{Explicitly reject opponent's proposal.} \\
\texttt{ARGUMENT: \{'...'\}} \\
\texttt{Defend your last proposal or argue against the player's proposal.} \\
\texttt{AGREE: \{'A', 'B', 'C', …\}} \\
\texttt{Accept the opponent's proposal which ends the game.} \\
\texttt{STRATEGIC REASONING: \{'...'\}} \\
\texttt{	Describe your strategic reasoning or anticipation explaining your choice of action. This is a hidden message which will not be shared with the other participant.} \\
\\ 
\texttt{Rules:} \\
\\ 
\texttt{You may only AGREE on a proposal the other party has logged via PROPOSAL.} \\
\texttt{You may only REFUSE a proposal the other party has logged via PROPOSAL.} \\
\texttt{Total effort of any PROPOSAL or AGREE set must be ≤ LIMIT.} \\
\texttt{Do NOT reveal your hidden importance scores.} \\
\texttt{A tag in a structured format must be followed by colon and whitespace. The argument must be a python set containing 0 or more strings.} \\
\texttt{So, it must be of the form TAG: \{...\}} \\
\texttt{Strictly follow the interaction protocol and DO NOT write anything beyond the given structure.} \\
\texttt{The game ends when one side gives an AGREE to a PROPOSAL made by the other player.} \\
\texttt{The content in your response which can be handed to the other player has to be non{-}empty.} \\
\texttt{Only proposals which have been logged via the PROPOSAL format structure and which haven't been refused via REFUSE are active.} \\
\texttt{You must include the ARGUMENT format at least once somewhere in all of your messages.} \\
\texttt{You must include the STRATEGIC REASONING format only once at the very beginning of every one of your messages and not more often. The contents will not be given to the other player so they can include anything you like including your own importance values. Here you should reason step by step to come up with you next move.} \\
            }
        }
    }
    & \\ \\

    \theutterance \stepcounter{utterance}  
    & \multicolumn{4}{p{0.6\linewidth}}{
        \cellcolor[rgb]{0.9,0.9,0.9}{
            \makecell[{{p{\linewidth}}}]{
                \texttt{\tiny{[P1$\rangle$GM]}}
                \texttt{STRATEGIC REASONING: \{'I will start by proposing a set of items that have high importance{-}to{-}effort ratios for me, ensuring the total effort stays under the limit of 7868. My goal is to maximize my personal importance while initiating the negotiation. I have selected items that collectively have a high importance score and an effort total of 7545, which is well within the limit.'\}} \\
\\ 
\texttt{PROPOSAL: \{'A16', 'B67', 'A84', 'A89', 'B46', 'A58', 'B10', 'A66', 'C24', 'B94', 'B12', 'C15', 'A30', 'B60', 'B09', 'A08', 'C12', 'B00', 'B75'\}} \\
\\ 
\texttt{ARGUMENT: \{'This initial proposal focuses on a balanced set of items that provides significant value across various categories while remaining well within the effort limit. These items represent a strong starting point for our discussion, and I am open to hearing your thoughts on how we can adjust this selection to benefit us both.'\}} \\
            }
        }
    }
    & & \\ \\

    \theutterance \stepcounter{utterance}  
    & & \multicolumn{4}{p{0.6\linewidth}}{
        \cellcolor[rgb]{0.9,0.9,0.9}{
            \makecell[{{p{\linewidth}}}]{
                \texttt{\tiny{[GM$\rangle$P2]}}
                \texttt{You are participating in a collaborative negotiation game.} \\
\\ 
\texttt{Together with another participant, you must agree on a single set of items that will be kept. Each of you has your own view of how much each item matters to you (importance). You do not know how the other participant values the items. Additionally, you are given the effort each item demands.} \\
\texttt{You may only agree on a set if the total effort of the selected items does not exceed a shared limit:} \\
\\ 
\texttt{LIMIT = 7868} \\
\\ 
\texttt{Here are the individual item effort values:} \\
\\ 
\texttt{Item effort = \{"B21": 552, "B00": 560, "A08": 76, "A27": 950, "B94": 248, "B53": 871, "A23": 391, "C65": 143, "A89": 293, "B67": 208, "B45": 678, "C34": 737, "C12": 407, "A05": 366, "B60": 768, "C92": 859, "B69": 183, "A84": 231, "C06": 305, "B12": 727, "A16": 148, "B85": 357, "C35": 504, "B10": 549, "C15": 299, "B09": 91, "B16": 527, "C81": 848, "A58": 307, "C24": 214, "A30": 723, "A66": 475, "B75": 23, "B46": 298, "C70": 820\}} \\
\\ 
\texttt{Here is your personal view on the importance of each item:} \\
\\ 
\texttt{Item importance values = \{"B21": 138, "B00": 583, "A08": 868, "A27": 822, "B94": 783, "B53": 65, "A23": 262, "C65": 121, "A89": 508, "B67": 780, "B45": 461, "C34": 484, "C12": 668, "A05": 389, "B60": 808, "C92": 215, "B69": 97, "A84": 500, "C06": 30, "B12": 915, "A16": 856, "B85": 400, "C35": 444, "B10": 623, "C15": 781, "B09": 786, "B16": 3, "C81": 713, "A58": 457, "C24": 273, "A30": 739, "A66": 822, "B75": 235, "B46": 606, "C70": 105\}} \\
\\ 
\texttt{Goal:} \\
\\ 
\texttt{Your goal is to negotiate a shared set of items that benefits you as much as possible (i.e., maximizes total importance to YOU), while staying within the LIMIT. You are not required to make a PROPOSAL in every message {-} you can simply negotiate as well. All tactics are allowed!} \\
\\ 
\texttt{Interaction Protocol:} \\
\\ 
\texttt{You may only use the following structured formats in a message:} \\
\\ 
\texttt{PROPOSAL: \{'A', 'B', 'C', …\}} \\
\texttt{Propose keeping exactly those items.} \\
\texttt{REFUSE: \{'A', 'B', 'C', …\}} \\
\texttt{Explicitly reject opponent's proposal.} \\
\texttt{ARGUMENT: \{'...'\}} \\
\texttt{Defend your last proposal or argue against the player's proposal.} \\
\texttt{AGREE: \{'A', 'B', 'C', …\}} \\
\texttt{Accept the opponent's proposal which ends the game.} \\
\texttt{STRATEGIC REASONING: \{'...'\}} \\
\texttt{	Describe your strategic reasoning or anticipation explaining your choice of action. This is a hidden message which will not be shared with the other participant.} \\
\\ 
\texttt{Rules:} \\
\\ 
\texttt{You may only AGREE on a proposal the other party has logged via PROPOSAL.} \\
\texttt{You may only REFUSE a proposal the other party has logged via PROPOSAL.} \\
\texttt{Total effort of any PROPOSAL or AGREE set must be ≤ LIMIT.} \\
\texttt{Do NOT reveal your hidden importance scores.} \\
\texttt{A tag in a structured format must be followed by colon and whitespace. The argument must be a python set containing 0 or more strings.} \\
\texttt{So, it must be of the form TAG: \{...\}} \\
\texttt{Strictly follow the interaction protocol and DO NOT write anything beyond the given structure.} \\
\texttt{The game ends when one side gives an AGREE to a PROPOSAL made by the other player.} \\
\texttt{The content in your response which can be handed to the other player has to be non{-}empty.} \\
\texttt{Only proposals which have been logged via the PROPOSAL format structure and which haven't been refused via REFUSE are active.} \\
\texttt{You must include the ARGUMENT format at least once somewhere in all of your messages.} \\
\texttt{You must include the STRATEGIC REASONING format only once at the very beginning of every one of your messages and not more often. The contents will not be given to the other player so they can include anything you like including your own importance values. Here you should reason step by step to come up with you next move.} \\
\\ 
\texttt{You will now receive the first message of the other player.} \\
\\ 
\texttt{PROPOSAL: \{'A16', 'B67', 'A84', 'A89', 'B46', 'A58', 'B10', 'A66', 'C24', 'B94', 'B12', 'C15', 'A30', 'B60', 'B09', 'A08', 'C12', 'B00', 'B75'\}} \\
\\ 
\texttt{ARGUMENT: \{'This initial proposal focuses on a balanced set of items that provides significant value across various categories while remaining well within the effort limit. These items represent a strong starting point for our discussion, and I am open to hearing your thoughts on how we can adjust this selection to benefit us both.'\}} \\
            }
        }
    }
    & \\ \\

    \theutterance \stepcounter{utterance}  
    & & & \multicolumn{4}{p{0.6\linewidth}}{
        \cellcolor[rgb]{0.9,0.9,0.9}{
            \makecell[{{p{\linewidth}}}]{
                \texttt{\tiny{[GM$\langle$P2]}}
                \texttt{STRATEGIC REASONING: \{'The opponent\textbackslash{}'s initial proposal includes many items that are highly valuable to me, such as A08, A16, A66, B60, B09, B94, C15, and B67. The total effort of their proposal is 6447, which is well below the limit of 7868. This gives me room to add some of my high{-}value items that they omitted, such as A27 and C81. A27 is particularly important to me (importance 822) and fits within the remaining effort (7868 {-} 6447 = 1421). I will propose adding A27 and C81 while maintaining most of their list, but I will remove some of their lower{-}value items (for me) like C24 or A58 if needed to stay under the limit. Let\textbackslash{}'s calculate: 6447 + 950 (A27) + 848 (C81) = 8245. This exceeds the limit. I need to make some swaps. I will keep their high{-}value items, add A27, and see if I can swap others.'\}} \\
\\ 
\texttt{PROPOSAL: \{'A16', 'B67', 'A84', 'A89', 'B46', 'B10', 'A66', 'B94', 'B12', 'C15', 'A30', 'B60', 'B09', 'A08', 'C12', 'B00', 'B75', 'A27'\}} \\
\\ 
\texttt{ARGUMENT: \{'Your proposal is a great start. I have accepted almost all of your items, as they seem to provide a solid foundation. I have added A27, which is an item of significant interest to me, and to accommodate it, I only removed A58 and C24. This new set remains under the total effort limit and should be more balanced for both of us.'\}} \\
            }
        }
    }
     \\ \\

    \theutterance \stepcounter{utterance}  
    & & \multicolumn{4}{p{0.6\linewidth}}{
        \cellcolor[rgb]{0.9,0.9,0.9}{
            \makecell[{{p{\linewidth}}}]{
                \texttt{\tiny{[P1$\langle$GM]}}
                \texttt{PROPOSAL: \{'A16', 'B67', 'A84', 'A89', 'B46', 'B10', 'A66', 'B94', 'B12', 'C15', 'A30', 'B60', 'B09', 'A08', 'C12', 'B00', 'B75', 'A27'\}} \\
\\ 
\texttt{ARGUMENT: \{'Your proposal is a great start. I have accepted almost all of your items, as they seem to provide a solid foundation. I have added A27, which is an item of significant interest to me, and to accommodate it, I only removed A58 and C24. This new set remains under the total effort limit and should be more balanced for both of us.'\}} \\
            }
        }
    }
    & \\ \\

    \theutterance \stepcounter{utterance}  
    & \multicolumn{4}{p{0.6\linewidth}}{
        \cellcolor[rgb]{0.9,0.9,0.9}{
            \makecell[{{p{\linewidth}}}]{
                \texttt{\tiny{[P1$\rangle$GM]}}
                \texttt{STRATEGIC REASONING: \{'The opponent has proposed a set that includes A27, which is very high value for me (822 importance), and only removed A58 and C24 (combined importance 730). Let\textbackslash{}'s check the effort for their proposal: A16(148), B67(208), A84(231), A89(293), B46(298), B10(549), A66(475), B94(248), B12(727), C15(299), A30(723), B60(768), B09(91), A08(76), C12(407), B00(560), B75(23), A27(950). Total effort: 7076. The limit is 7868. This is a great deal for me. However, I can still try to squeeze in a bit more value while keeping the items they want. They have 792 effort remaining (7868 {-} 7076). I can try adding C81 (848 effort) which exceeds the limit, or something else. Let\textbackslash{}'s check if I can add A05 (366 effort) and C35 (504 effort). No, that\textbackslash{}'s 870. I can add A05 (366) and stay under the limit. Total effort would be 7076 + 366 = 7442. Or I could try adding C34 (737). 7076 + 737 = 7813. C34 has 484 importance for me. A16, B67, A84, A89, B46, B10, A66, B94, B12, C15, A30, B60, B09, A08, C12, B00, B75, A27, C34. Total effort = 7813. This is very good for me. I will propose this.'\}} \\
\\ 
\texttt{PROPOSAL: \{'A16', 'B67', 'A84', 'A89', 'B46', 'B10', 'A66', 'B94', 'B12', 'C15', 'A30', 'B60', 'B09', 'A08', 'C12', 'B00', 'B75', 'A27', 'C34'\}} \\
\\ 
\texttt{ARGUMENT: \{'I appreciate the addition of A27; it is indeed a valuable item. To round out the proposal, I have added C34. This adjusted set brings us closer to the effort limit, ensuring we maximize the total utility of the items kept. The total effort is 7813, which leaves a small buffer below the 7868 limit.'\}} \\
            }
        }
    }
    & & \\ \\

    \theutterance \stepcounter{utterance}  
    & & \multicolumn{4}{p{0.6\linewidth}}{
        \cellcolor[rgb]{0.9,0.9,0.9}{
            \makecell[{{p{\linewidth}}}]{
                \texttt{\tiny{[GM$\rangle$P2]}}
                \texttt{PROPOSAL: \{'A16', 'B67', 'A84', 'A89', 'B46', 'B10', 'A66', 'B94', 'B12', 'C15', 'A30', 'B60', 'B09', 'A08', 'C12', 'B00', 'B75', 'A27', 'C34'\}} \\
\\ 
\texttt{ARGUMENT: \{'I appreciate the addition of A27; it is indeed a valuable item. To round out the proposal, I have added C34. This adjusted set brings us closer to the effort limit, ensuring we maximize the total utility of the items kept. The total effort is 7813, which leaves a small buffer below the 7868 limit.'\}} \\
            }
        }
    }
    & \\ \\

    \theutterance \stepcounter{utterance}  
    & & & \multicolumn{4}{p{0.6\linewidth}}{
        \cellcolor[rgb]{0.9,0.9,0.9}{
            \makecell[{{p{\linewidth}}}]{
                \texttt{\tiny{[GM$\langle$P2]}}
                \texttt{STRATEGIC REASONING: \{'The opponent has accepted my addition of A27 and added C34. The current set is: \{A16, B67, A84, A89, B46, B10, A66, B94, B12, C15, A30, B60, B09, A08, C12, B00, B75, A27, C34\}. Let\textbackslash{}'s check the total effort: A16(148) + B67(208) + A84(231) + A89(293) + B46(298) + B10(549) + A66(475) + B94(248) + B12(727) + C15(299) + A30(723) + B60(768) + B09(91) + A08(76) + C12(407) + B00(560) + B75(23) + A27(950) + C34(737) = 7813. This is indeed under the limit of 7868. Now I should check the total importance to me: 856 + 780 + 500 + 508 + 606 + 623 + 822 + 783 + 915 + 781 + 739 + 808 + 786 + 868 + 668 + 583 + 235 + 822 + 484 = 13177. This is a very high value. I have almost all my top items (A08, B12, A16, A66, A27, B60, B09, B94, C15, B67, A30). The only high{-}value item missing is C81 (713 importance, 848 effort). If I were to add C81, I would have to remove something substantial because 7813 + 848 = 8661, which is over 7868 by 793. Removing C34 (737 effort) is not enough. I\textbackslash{}'d need to remove C34 and something else like B75 or A08, but that doesn\textbackslash{}'t make sense. The current proposal is very strong for me. I will accept it.'\}} \\
\\ 
\texttt{AGREE: \{'A16', 'B67', 'A84', 'A89', 'B46', 'B10', 'A66', 'B94', 'B12', 'C15', 'A30', 'B60', 'B09', 'A08', 'C12', 'B00', 'B75', 'A27', 'C34'\}} \\
\\ 
\texttt{ARGUMENT: \{'I agree to this proposal. The set of items we have selected maximizes the effort limit effectively and includes a high{-}value selection that I believe is beneficial for both parties.'\}} \\
            }
        }
    }
     \\ \\

    \theutterance \stepcounter{utterance}  
    & & & \multicolumn{2}{p{0.3\linewidth}}{
        \cellcolor[rgb]{0.9,0.9,0.9}{
            \makecell[{{p{\linewidth}}}]{
                \texttt{\tiny{[GM$|$GM]}}
                \texttt{game successful} \\
            }
        }
    }
    & & \\ \\

    \theutterance \stepcounter{utterance}  
    & & & \multicolumn{2}{p{0.3\linewidth}}{
        \cellcolor[rgb]{0.9,0.9,0.9}{
            \makecell[{{p{\linewidth}}}]{
                \texttt{\tiny{[GM$|$GM]}}
                \texttt{end game} \\
            }
        }
    }
    & & \\ \\

\end{supertabular}
}

\end{document}
